% Rebuttal to Reviewer Comments
% ACM-BCB 2026 Submission: Structure-Aware Prediction of PROTAC-Mediated
% Protein Degradability via Graph Neural Networks
%
% Compile: pdflatex rebuttal.tex && pdflatex rebuttal.tex

\documentclass[11pt]{article}
\usepackage[margin=1in]{geometry}
\usepackage{xcolor}
\usepackage{booktabs}
\usepackage{enumitem}
\usepackage{hyperref}
\usepackage{amsmath}

% Styling
\definecolor{revblue}{RGB}{0,70,170}
\newcommand{\reviewerbox}[2]{%
  \noindent\fcolorbox{black}{gray!10}{\parbox{\dimexpr\linewidth-2\fboxsep-2\fboxrule}{%
    \textbf{#1}: #2}}%
  \vspace{4pt}%
}
\newcommand{\resp}{\noindent\textbf{Response:}~}
\newcommand{\change}[1]{\textcolor{revblue}{[#1]}}

\setlength{\parskip}{6pt}
\setlength{\parindent}{0pt}

\begin{document}

\begin{center}
{\Large\bfseries Response to Reviewer Comments}\\[6pt]
{\large Structure-Aware Prediction of PROTAC-Mediated Protein Degradability\\via Graph Neural Networks}\\[6pt]
ACM-BCB 2026 Submission
\end{center}

\bigskip

We thank all three reviewers for their thorough and constructive feedback.
We have substantially revised the manuscript, and all changes are
\textcolor{revblue}{highlighted in blue} in the revised paper.
Line numbers below refer to the revised \texttt{paper.tex}.
The major revisions include:

\begin{itemize}[topsep=2pt,itemsep=2pt]
  \item \textbf{6-seed extended validation} (seeds 42, 123, 456, 789, 1011, 1213): the 6-seed mean is 0.603$\pm$0.097, honestly reported alongside the original 3-seed mean (0.646$\pm$0.124).
  \item \textbf{Statistical significance testing}: bootstrap $p$-values confirm the improvement over gradient boosting is not significant ($p=0.259$ for 3-seed, $p=0.556$ for 6-seed). This is stated explicitly in the main text.
  \item \textbf{VHL root cause analysis}: we rule out class imbalance and dataset size, identifying structural heterogeneity as the likely cause of VHL failure (0.396 AUROC).
  \item \textbf{BRD4 biological case study}: lysine attention weights identify K374 as the top-weighted residue, validated by mass spectrometry literature.
  \item \textbf{Calibration analysis}: ECE\,=\,0.029 on target-unseen (excellent), providing a new quantitative strength.
  \item \textbf{Narrowed claims}: E3 generalization narrowed to CRBN$\leftrightarrow$VHL transfer; multi-seed average is the primary reported result; limitations stated prominently before results.
\end{itemize}

\bigskip
\hrule
\bigskip

%% ====================================================================
%%  REVIEWER 9v2E
%% ====================================================================

\section*{Response to Reviewer 9v2E (Rating: 4 --- Borderline Accept)}

\reviewerbox{C1}{The dataset is relatively small (3,101 samples, 155 unique proteins, 10 E3 ligases). The target-unseen effective diversity is much closer to the number of unique proteins, contributing to seed sensitivity.}

\resp We fully agree. We have added a prominent ``Dataset characteristics and limitations'' paragraph \emph{before} any results are presented:

\begin{quote}\itshape
``PROTAC-8K contains 3,101 labeled samples spanning only 155 unique protein targets\ldots\ The E3 ligase distribution is highly imbalanced: CRBN accounts for 62\% of samples, VHL for 35\%, and the remaining eight E3s combined represent only 3\%.\ These constraints fundamentally bound generalization performance\ldots''
\end{quote}

We also note that PROTAC-8K is the only publicly available benchmark of this scale, and we have no external validation dataset---both stated as explicit limitations.

\change{Lines 444--449 (new paragraph before results); Limitations items 1, 2, 6 at lines 808, 809, 813.}

\bigskip

\reviewerbox{C2}{Structural validation is limited. Predictions are not supported by docking, ternary-complex modeling, MD, or free energy calculations.}

\resp We agree this is a limitation. It is now discussed in the Limitations paragraph:

\begin{quote}\itshape
``Predictions rely on static AlphaFold structures without physics-based validation (docking, MD, free energy calculations). The BRD4 case study (K374) provides one anecdotal positive control.''
\end{quote}

We note an inherent tension: ternary-complex modeling requires knowing the PROTAC linker geometry, which is unavailable in our pre-synthesis setting. Post-hoc structural validation against experimentally determined ternary complexes (where available) is an important future direction.

\change{Limitation item~5 at line 812.}

\bigskip

\reviewerbox{C3}{No compelling biological case study showing why a specific protein is predicted degradable, which lysines drive the prediction, or how predictions align with literature.}

\resp We have added a new subsection ``Case Study: BRD4 Degradability'' that provides exactly this analysis:

\begin{itemize}[topsep=2pt,itemsep=1pt]
  \item BRD4 (O60885) is predicted highly degradable with CRBN (score: 0.87).
  \item Lysine attention weights identify \textbf{K374} (attention: 0.12, SASA: 94\,\AA$^2$) as the top residue---located in a disordered, solvent-exposed linker between bromodomains.
  \item Literature validation: mass spectrometry studies by Zengerle et al.\ (2015) confirm K374 as a major ubiquitination site following JQ1-based PROTAC treatment.
\end{itemize}

We acknowledge this is a single successful case and that systematic validation across more targets would strengthen confidence.

\change{Lines 719--736 (new \S4.5 ``Case Study: BRD4 Degradability'').}

\bigskip

\reviewerbox{Additional}{The key target-unseen result is unstable; the paper emphasizes the best seed over the multi-seed result; some biological and E3 generalization claims are only partly supported.}

\resp All three points are now addressed:

\begin{enumerate}[topsep=2pt,itemsep=1pt]
  \item \textbf{Instability}: 6-seed extended validation (0.603$\pm$0.097) reported alongside the original 3-seed (0.646$\pm$0.124); ensembling across $\geq$6 seeds is explicitly recommended for deployment (lines 651--656).
  \item \textbf{Best seed de-emphasized}: the multi-seed average is the primary result throughout; best seed (0.7449) is consistently marked as secondary, parenthetical (abstract line 65; Table~1 at line 457; conclusion line 825).
  \item \textbf{E3 claims narrowed}: changed throughout from ``E3 generalization'' to ``CRBN$\to$VHL transfer'' (abstract line 65; section heading line 667; conclusion line 825); rare E3 generalization explicitly stated as unsupported (line 680).
\end{enumerate}

\change{Abstract lines 65--69; Introduction lines 135--138; \S4.3 lines 651--680; Discussion lines 780--784; Conclusion lines 825--827.}


\bigskip
\hrule
\bigskip

%% ====================================================================
%%  REVIEWER Y2rf
%% ====================================================================

\section*{Response to Reviewer Y2rf (Rating: 2 --- Reject)}

We appreciate the reviewer's detailed and rigorous critique. We have addressed every major and minor concern and believe the revised manuscript warrants reconsideration.

\bigskip

\reviewerbox{M1}{The main result should be multi-seed average, not best seed. Rewrite abstract, results, and conclusion accordingly.}

\resp Done. Throughout the revised manuscript:

\begin{itemize}[topsep=2pt,itemsep=1pt]
  \item The abstract leads with ``0.646$\pm$0.124 AUROC on target-unseen evaluation (best seed: 0.7449)'' (line 65).
  \item The introduction reports both 3-seed and 6-seed means as primary numbers (lines 135--136).
  \item Table~1 (line 457) now reports both 3-seed (0.646) and 6-seed (0.603) means, with ``Best Seed'' as a separate column.
  \item Section 4.2 states: ``We report multi-seed averages as primary results\ldots\ best-seed performance (0.7449) represents peak achievable performance but should not be interpreted as typical expectation'' (lines 471--474).
  \item The conclusion leads with average performance (line 825).
\end{itemize}

\change{Lines 65, 135--136, 457--469 (Table~1), 471--474, 825--827.}

\bigskip

\reviewerbox{M2}{Target-unseen performance is unstable. Report more seeds and provide confidence intervals. Discuss average performance rather than peak.}

\resp We trained 3 additional seeds (789, 1011, 1213), yielding AUROCs of 0.657, 0.520, and 0.502 respectively. The 6-seed results are:

\begin{center}
\begin{tabular}{lcc}
\toprule
\textbf{Estimate} & \textbf{AUROC} & \textbf{95\% CI} \\
\midrule
3-seed mean & 0.646 $\pm$ 0.124 & [0.506, 0.745] \\
6-seed mean & 0.603 $\pm$ 0.097 & [0.532, 0.682] \\
Gradient Boosting & 0.607 & --- \\
\bottomrule
\end{tabular}
\end{center}

The 6-seed mean (0.603) is marginally \emph{below} the gradient boosting baseline (0.607, bootstrap $p=0.556$). We report this honestly and prominently:

\begin{quote}\itshape
``This result is important to report honestly: the original 3-seed mean (0.646) was moderately favorable; the 6-seed analysis demonstrates that DegradoMap's advantage over hand-crafted feature baselines is not robust across all initializations.'' (lines 653--654)
\end{quote}

We explicitly recommend ensembling across $\geq$6 seeds for deployment (line 656) and note that individual seeds span a 0.25-AUROC range.

\change{Lines 651--657 (``Six-seed extended validation'' paragraph); Table~1 lines 457--469; Figure~4 at line 664.}

\bigskip

\reviewerbox{M3}{The lysine-related biological claim is not supported by ablation---removing the lysine indicator \emph{improves} performance. Explain or retract.}

\resp We now discuss this ``lysine indicator paradox'' directly in the Discussion with concrete evidence:

\begin{itemize}[topsep=2pt,itemsep=1pt]
  \item With lysine indicator: val 0.584, test 0.477 (val-test gap: $-$0.106 $\to$ \textbf{overfitting}).
  \item Without lysine indicator: val 0.541, test 0.578 (val-test gap: $+$0.038 $\to$ generalizes).
\end{itemize}

This is a classic overfitting pattern: the lysine indicator improves validation but hurts test performance. We propose that the lysine-weighted \emph{pooling mechanism} (Eq.~3) already encodes lysine information architecturally, making the binary node feature redundant and causing overfitting to training-set lysine distributions (lines 791--795).

We conclude: ``architectural inductive biases (lysine pooling) are more effective than explicit feature annotation for this task,'' and note that current feature design is suboptimal (Limitation~7, line 814).

\change{Lines 789--796 (``Lysine indicator paradox'' paragraph); Figure~7 at line 804; Limitation~7 at line 814.}

\bigskip

\reviewerbox{M4}{The E3 generalization claim should be narrower. Evidence supports mainly CRBN$\leftrightarrow$VHL transfer.}

\resp Agreed and revised throughout:

\begin{itemize}[topsep=2pt,itemsep=1pt]
  \item Abstract: ``CRBN$\to$VHL E3-unseen transfer'' (line 65).
  \item Section heading: ``E3 generalization is limited to major ligases'' (line 667).
  \item New per-E3 Table~5 (line 686) and Figure~5 (line 700) show: CRBN 0.758, VHL 0.396, cIAP1 0.417.
  \item VHL root cause analysis (lines 672--680): ruled out class imbalance (36.9\% vs 38.4\%) and dataset size (91 vs 125 proteins); identified structural heterogeneity as the likely cause.
  \item Conclusion: ``CRBN$\to$VHL transfer (0.811 AUROC)'' (line 825).
\end{itemize}

\change{Lines 65, 667--680, 686--696 (Table~5), 700--704 (Figure~5), 825.}

\bigskip

\reviewerbox{M5}{Validation setup should better align with the target-unseen setting, or explain why the current setup is sufficient.}

\resp We now explicitly acknowledge this as a methodological limitation:

\begin{quote}\itshape
``Validation-split AUROC does not reliably predict target-unseen test AUROC\ldots\ Ideally, we would use a target-unseen validation split for hyperparameter selection; however, this would further reduce the already-limited training data (155 proteins $\to$ ${\sim}$120 training, ${\sim}$15 validation, ${\sim}$20 test).'' (lines 340--341)
\end{quote}

We state that the chosen hyperparameters may not be optimal for target-unseen generalization, and that this likely contributes to variance (line 343). This limitation is also listed as item~4 in the Limitations paragraph (line 811).

\change{Lines 339--343 (``Validation-test mismatch'' paragraph in \S3.5); Limitation~4 at line 811.}

\bigskip

\reviewerbox{Minor 1}{Dataset limitations should be stated more directly in main results.}

\resp A ``Dataset characteristics and limitations'' paragraph now appears at the start of Results, \emph{before} any performance numbers (see response to Reviewer 9v2E, C1).

\change{Lines 444--449.}

\bigskip

\reviewerbox{Minor 2}{The lack of external validation is an important limitation.}

\resp Now stated as Limitation~6:

\begin{quote}\itshape
``No external validation: PROTAC-8K is the only publicly available benchmark with sufficient scale; generalization to data from other assay platforms or research groups is unvalidated.'' (line 813)
\end{quote}

\change{Line 813.}


\bigskip
\hrule
\bigskip

%% ====================================================================
%%  REVIEWER PtZf
%% ====================================================================

\section*{Response to Reviewer PtZf (Rating: 3 --- Borderline Reject)}

\reviewerbox{W1}{Performance is inconsistent across the paper: abstract reports 0.78 AUROC (95\% CI 0.74--0.82), main results section reports 0.646$\pm$0.124, appendix reports 5-fold CV mean 0.565.}

\resp We believe the ``0.78'' figure may refer to a prior draft; the current abstract reports 0.646$\pm$0.124 (multi-seed average, line 65). To prevent any confusion, we have added a ``Reconciling performance estimates'' paragraph that explicitly maps each reported number to its evaluation protocol:

\begin{center}
\begin{tabular}{lll}
\toprule
\textbf{Number} & \textbf{Protocol} & \textbf{Role} \\
\midrule
0.603 $\pm$ 0.097 & 6-seed, single split & Most conservative \\
0.646 $\pm$ 0.124 & 3-seed, single split & Primary result \\
0.7449 & Best seed & Peak (secondary) \\
0.657 & Baseline (seed 42) & Single-seed reference \\
0.565 $\pm$ 0.052 & 5-fold CV & Conservative lower bound \\
\bottomrule
\end{tabular}
\end{center}

The 6-seed mean is recommended for citation as the most conservative estimate.

\change{Lines 712--716 (``Reconciling performance estimates'' paragraph in \S4.3).}

\bigskip

\reviewerbox{W2}{Robustness is limited. Target-unseen has high seed sensitivity, and 5-fold CV mean is substantially lower than the headline single-split number.}

\resp We agree and have:
\begin{enumerate}[topsep=2pt,itemsep=1pt]
  \item Extended to 6 seeds, yielding 0.603$\pm$0.097 --- marginally below gradient boosting (0.607); lines 651--657.
  \item Stated explicitly that the improvement is \textbf{not statistically significant}: ``bootstrap p = 0.259'' (line 550) and ``p = 0.556'' (line 653).
  \item Recommended ensembling across $\geq$6 seeds for reliable deployment (line 656).
  \item Reported both 3-seed and 6-seed results in Table~1 (lines 457--469) to show how the estimate changes with more seeds.
\end{enumerate}

\change{Lines 457--469 (Table~1), 550--552, 651--657, 780--784 (``Honest performance assessment'').}

\bigskip

\reviewerbox{W3}{Validation and model-selection protocol appears weak for the main use case. Validation AUROC does not reliably predict target-unseen test AUROC.}

\resp Acknowledged as a methodological limitation (see response to Reviewer Y2rf, M5). We now discuss this explicitly in the Training Procedure section and Limitation~4.

\change{Lines 339--343 (``Validation-test mismatch''); Limitation~4 at line 811.}

\bigskip

\reviewerbox{W4}{Generalization across E3 ligases remains limited. VHL performance is poor; rare ligases have too few samples.}

\resp Fully addressed (see response to Reviewer Y2rf, M4). E3 claims are narrowed to CRBN$\leftrightarrow$VHL transfer. A new per-E3 table (Table~5, line 686) and VHL root cause analysis (lines 672--680) are included. We now explicitly state:
\begin{quote}\itshape
``practitioners should expect substantially lower performance than the overall 0.646 mean when deploying on VHL-targeted proteins, and near-random performance for rare E3s.'' (lines 679--680)
\end{quote}

\change{Lines 667--680, 686--696 (Table~5), 700--704 (Figure~5).}

\bigskip

\reviewerbox{W5}{No external validation reduces confidence in real-world generalization.}

\resp Stated as Limitation~6 (see response to Reviewer Y2rf, Minor~2). We caution against deployment without prospective validation on independently collected data.

\change{Line 813.}

\bigskip
\hrule
\bigskip

\section*{Summary of All Changes}

The table below maps each reviewer concern to the specific lines changed in the revised manuscript.

\begin{center}
\small
\begin{tabular}{@{}p{3.8cm}p{2.2cm}p{6.5cm}@{}}
\toprule
\textbf{Concern} & \textbf{Reviewer(s)} & \textbf{Revision (lines in \texttt{paper.tex})} \\
\midrule
Best seed over-emphasized & Y2rf M1, PtZf W1 & Abstract L65; Table~1 L457--469; \S4.2 L471--474; Conclusion L825--827 \\
Target-unseen unstable & Y2rf M2, PtZf W2 & \S4.3 L651--657; Table~1 L463; Figure~4 L664; $p$-values L474, L550, L653 \\
Lysine claim unsupported & Y2rf M3 & \S5 L789--796; Figure~7 L804; Limitation~7 L814 \\
E3 claims too broad & Y2rf M4, PtZf W4 & Abstract L65; \S4.3 L667--680; Table~5 L686; Figure~5 L700; Conclusion L825 \\
Validation misaligned & Y2rf M5, PtZf W3 & \S3.5 L339--343; Limitation~4 L811 \\
Dataset limits buried & 9v2E C1, Y2rf M6 & \S4.1 L444--449 \\
No structural validation & 9v2E C2 & Limitation~5 L812 \\
No case study & 9v2E C3 & \S4.5 L719--736 \\
Inconsistent reporting & PtZf W1 & \S4.3 L712--716 (reconciliation paragraph) \\
No external validation & Y2rf M7, PtZf W5 & Limitation~6 L813 \\
\midrule
\multicolumn{3}{@{}l}{\textit{New analyses added in revision:}} \\
\midrule
6-seed validation & --- & Table~1 L463; \S4.3 L651--657; Figure~4 L664 \\
Calibration analysis & --- & \S4.6 L738--765; Table~6 L747 \\
Per-E3 breakdown & --- & Table~5 L686; Figure~5 L700 \\
VHL root cause & --- & \S4.3 L672--680 \\
BRD4 case study & --- & \S4.5 L719--736 \\
Lysine paradox figure & --- & Figure~7 L804 \\
Statistical significance & --- & L474, L550, L653 \\
Honest assessment & --- & \S5 L780--784 \\
\bottomrule
\end{tabular}
\end{center}

\end{document}
